\chapter{Discusión}

\section{Búsqueda de hiperparámetros}

El HV medio alcanzado por las mejores configuraciones depende principalmente de la dificultad del problema de clasificación y no de la dimensionalidad de entrada (aproximadamente 0.93 en \texttt{ds2}, 0.87 en \texttt{ds1}, 0.81 en \texttt{ds4} y 0.75 en \texttt{ds3}, sin relación monótona con el número de atributos), algo esperable dado que el HV combina simultáneamente exactitud y complejidad estructural. El análisis de importancia de hiperparámetros muestra que la escala de pesos (\texttt{wScale}) es sistemáticamente el hiperparámetro más influyente en tres de los cuatro estudios, coherente con su papel de amplificador necesario para que las conexiones codificadas entre [0,1] generen disparos efectivos en la red; las excepciones son \texttt{ds2}, dominado por la variable \texttt{architecture}, y \texttt{ds4}, donde la importancia se reparte entre varios hiperparámetros sin que ninguno domine, en un espacio de mayor dimensionalidad. Bajo un mismo presupuesto de búsqueda, MOEA-CKF y S-NSGA-II alcanzan HV prácticamente equivalentes en la configuración ganadora: en tres de los cuatro datasets (\texttt{ds1}, \texttt{ds3} y \texttt{ds4}) S-NSGA-II obtiene un HV medio ligeramente superior, con diferencias entre 0.002 y 0.016, mientras que en \texttt{ds2} la relación se invierte levemente a favor de MOEA-CKF, con una diferencia de apenas 0.0003; esto sugiere que, bajo el mismo presupuesto, ambos SMOEA podrían alcanzar HV comparables, aunque esta lectura se basa en el mejor \emph{trial} de un único estudio y se contrasta con mayor rigor estadístico en el Estudio 1. Estos resultados justifican tratar \texttt{architecture}, \texttt{wScale} y los límites de esparsidad como hiperparámetros externos del pipeline y sustentan las configuraciones utilizadas en la fase de comparación de algoritmos.

\section{Estudio 1: Comparación de algoritmos en SNN}
\label{sec:discusion-estudio1}

\subsection{Hipervolumen}

En hipervolumen (HV), S-NSGA-II obtiene una ventaja estadísticamente significativa en dos de los cuatro datasets (\texttt{ds1}, \texttt{ds3}), MOEA-CKF en uno (\texttt{ds4}), y en el restante (\texttt{ds2}) no se observa diferencia significativa. La magnitud de esta ventaja varía sustancialmente entre datasets: es más marcada en el de mayor dimensionalidad, \texttt{ds4} —donde además se invierte a favor de MOEA-CKF—, y es prácticamente inexistente en \texttt{ds2} (18 características), donde ambos algoritmos convergen a una región estrecha y muy alta del espacio de soluciones (diferencia mediana de solo 0.00016 entre medianas, con márgenes mínimos en las corridas donde gana S-NSGA-II y márgenes mayores en las pocas donde gana MOEA-CKF). Estos patrones son consistentes, aunque no concluyentes con solo cuatro puntos de datos, con la hipótesis de que S-NSGA-II rinde mejor en problemas de menor dimensionalidad (14 y 24 características) y MOEA-CKF en el de mayor dimensionalidad (60 características). Adicionalmente, en tres de los cuatro datasets (\texttt{ds1}, \texttt{ds2}, \texttt{ds3}) MOEA-CKF muestra menor varianza entre corridas que S-NSGA-II, incluso cuando pierde en mediana: tiende a ser más consistente y predecible entre ejecuciones, una propiedad relevante si el criterio de evaluación valora la reproducibilidad además del desempeño puntual.

\subsection{Frentes de Pareto}

El patrón de complejidad-error observado en los frentes explica en gran medida las diferencias de HV descritas arriba: en los datasets donde gana S-NSGA-II (\texttt{ds1}, \texttt{ds3}), sus redes de mejor accuracy son sistemáticamente más densas y más exactas que las de MOEA-CKF, mientras que en \texttt{ds4}, donde gana MOEA-CKF, los roles se invierten y es este último quien produce redes más densas y más precisas. La cobertura mutua de los frentes combinados (C-metric) refuerza esta lectura: es relativamente pareja en \texttt{ds2} y \texttt{ds3} (en torno a 0.50/0.50 y 0.31/0.33) y marcadamente asimétrica en \texttt{ds4} (0.86/0.14 a favor de MOEA-CKF), de modo que el dataset con la ventaja de HV más contundente es también el de mayor dominancia punto a punto.

Sin embargo, HV y dominancia punto a punto no siempre coinciden: en \texttt{ds1}, pese a tener menos puntos en su frente combinado, el frente de MOEA-CKF domina en mayor proporción al de S-NSGA-II que al revés (55.6\% frente a 16.7\%), lo que indica que la ventaja de HV de S-NSGA-II proviene de cubrir una región más amplia del espacio de soluciones —en particular la banda de mayor complejidad, donde concentra la mayoría de sus puntos— y no de que cada solución individual sea mejor que su contraparte. Un mecanismo relacionado aparece en \texttt{ds3}, donde S-NSGA-II traduce el incremento de complejidad estructural en una reducción más sostenida del error a partir de cierto umbral, lo que explica que su frente combinado sea sustancialmente más grande y contribuye a su ventaja de HV. En \texttt{ds2}, en cambio, S-NSGA-II parece especializarse en soluciones muy esparsas con buenos resultados, mientras que MOEA-CKF explora una región más amplia sin obtener una ganancia proporcional, lo cual es coherente con que el HV no distinga a los algoritmos en ese dataset.

Dos limitaciones metodológicas condicionan esta lectura. Primero, el análisis de frentes se apoya en la solución de menor error de entrenamiento de cada corrida como sustituto de la accuracy de test, ya que el 100\% de las corridas archivadas usan arquitecturas de dos salidas (WTA) y no es posible medir la generalización de forma directa; se trata de una limitación estructural, no de un dato pendiente de calcular, y el análisis complejidad-error de todo el frente es el sustituto disponible. Segundo, el C-metric sobre los frentes combinados debe interpretarse como una magnitud aproximada y no como una proporción de alta precisión: tras deduplicar y filtrar soluciones no dominadas, estos frentes quedan reducidos a entre 4 y 14 puntos, lo que limita la resolución de la métrica en conjuntos tan pequeños.

\subsection{Convergencia}

En los tres datasets donde existe una diferencia clara de HV (\texttt{ds1}, \texttt{ds3}, \texttt{ds4}), el algoritmo con menor HV final converge sistemáticamente más rápido a su propio techo, mientras que el de mayor HV final asciende más lentamente hacia uno más alto: MOEA-CKF llega antes a un techo más bajo en \texttt{ds1} y \texttt{ds3} ($\sim4.3\times$ más rápido en \texttt{ds1}), mientras que en \texttt{ds4} es S-NSGA-II quien se estabiliza antes en un nivel inferior y MOEA-CKF quien tarda $\sim2.4\times$ más en alcanzar el suyo, más alto, sin llegar a estabilizarse del todo dentro del presupuesto de evaluaciones. Un \textit{plateau} rápido en un techo bajo frente a un ascenso lento hacia un techo alto parece ser, entonces, un balance recurrente en este estudio, y no el ruido de una corrida aislada. En \texttt{ds2}, la ausencia de una diferencia relevante de HV se refleja también en la convergencia: la población inicial ya alcanza el 90\% del HV final para ambos algoritmos, por lo que la velocidad de convergencia mide más bien qué tan rápido se estabiliza la búsqueda que una carrera real de optimización sobre un problema comparativamente saturado.

Una limitación transversal a este análisis es que el eje de evaluaciones (FE) de las curvas y umbrales de velocidad es una aproximación por interpolación lineal dentro de cada corrida: no se registra el FE por fila sino solo la \texttt{generation}, que no es directamente comparable entre algoritmos porque MOEA-CKF consume más FE por generación que S-NSGA-II.

\section{Estudio 2: Plataforma SNN (C++) vs.\ ANN (MATLAB)}
\label{sec:discusion-estudio2}

Es el hallazgo más uniforme de todo el análisis: la implementación ANN original de PlatEMO/MATLAB supera a la SNN en HV y exactitud en las ocho combinaciones de dataset y algoritmo evaluadas, siempre de forma significativa tras corrección. La brecha de exactitud es menor en \texttt{ds1} y mayor en \texttt{ds4} para S-NSGA-II. La brecha de complejidad muestra, además, un patrón dependiente del algoritmo que se repite entre datasets: para S-NSGA-II, las redes ANN de mejor exactitud tienden a ser casi completamente densas frente a redes SNN mucho más esparsas (particularmente marcado en \texttt{ds2} y \texttt{ds3}), mientras que para MOEA-CKF esta brecha es sistemáticamente más moderada. La única excepción a este patrón ocurre en \texttt{ds4}, donde es la red SNN de MOEA-CKF la que resulta más densa que la ANN, la única anomalía de complejidad entre las ocho combinaciones plataforma-algoritmo del estudio. En conjunto, estos resultados sugieren que la SNN podría favorecer soluciones más esparsas que su contraparte ANN para una parte importante de las combinaciones algoritmo-dataset, un patrón dependiente del algoritmo que podría merecer investigación aparte.

MATLAB/PlatEMO original y la migración a C++ difieren en más que el modelo de neurona, ya que también difieren en la implementación y en el generador de números aleatorios. Por ello, los resultados de esta comparación deben leerse como evidencia de que la implementación ANN concreta supera consistentemente a la implementación SNN concreta evaluada aquí, y no como una conclusión general sobre la neurona de disparo frente a una neurona de tasa.

\section{Otros resultados}
\label{sec:discusion-otros}

El efecto del remuestreo sobre el HV final es heterogéneo entre datasets y algoritmos, sin que se identifique una regla simple. Mejora el HV de forma clara y con efecto grande en ambos algoritmos en \texttt{ds1} y \texttt{ds2} —este último con el menor costo relativo de tiempo de los cuatro datasets ($\times4.8$--$4.9$, frente a $\times5$--$15$ en el resto)—, pero produce resultados mixtos en \texttt{ds3} y \texttt{ds4}: en \texttt{ds3} solo mejora a MOEA-CKF, mientras que en S-NSGA-II no hay efecto significativo ($p=0.066$), con la dirección del tamaño de efecto incluso sugiriendo que la evaluación normal iguala o supera al remuestreo más seguido que a la inversa; en \texttt{ds4} el patrón se invierte y es solo S-NSGA-II quien mejora significativamente. El costo de tiempo del remuestreo, en cambio, es siempre sustancial (entre $\times4.8$ y $\times15.1$) independientemente de si el HV mejora o no, lo que deja como pregunta abierta si existe alguna interacción entre el remuestreo y los mecanismos específicos de cada algoritmo que explique esta heterogeneidad.

En tiempo de cómputo, S-NSGA-II es significativamente más rápido que MOEA-CKF en los cuatro datasets, con efecto grande y menor variabilidad entre corridas en todos los casos. La brecha, sin embargo, crece marcadamente con la dimensionalidad del problema: de aproximadamente $6$--$10\%$ en \texttt{ds1}--\texttt{ds3} (14--24 características) a cerca de $2.4\times$ en \texttt{ds4} (60 características, sin una sola excepción en las 15 corridas pareadas), donde MOEA-CKF resulta además el más lento y variable. Este patrón es consistente con que el mecanismo de fusión de conocimiento entre escalas de MOEA-CKF (SVD, análisis de dispersión) consume más FE por generación y escala peor con la dimensionalidad que los operadores más simples de S-NSGA-II.

Estas comparaciones de tiempo solo son válidas dentro del conjunto experimental temporal, ejecutado con hiperparámetros compartidos entre ambos algoritmos; los conjuntos experimentales normal y de remuestreo no permiten una comparación válida en tiempo de ejecución debido al paralelismo y los hilos no controlados entre tandas de ejecución, por lo que sus conclusiones de tiempo no deben mezclarse con las conclusiones de calidad (HV) del resto del análisis, que sí emplea hiperparámetros afinados.

\section{Conectividad de las redes}
\label{sec:discusion-conectividad}

El patrón de neuronas muertas y clases de salida desconectadas descrito en los resultados agregados no es producto del azar, sino una consecuencia directa del diseño de los operadores de dispersión de ambos algoritmos. Ninguno de los mecanismos que apagan pesos durante la optimización posee noción alguna de a qué neurona o a qué salida pertenece cada sinapsis: todos operan gen por gen, o sobre un conteo u objetivo global de dispersión, sin considerar en ningún momento la estructura del grafo subyacente. En la inicialización de la red, cada gen se apaga de forma independiente con una probabilidad aproximada del 50\%, un procedimiento puramente bernoulliano que no distingue entre los pesos que alimentan una neurona oculta y los que alimentan una neurona de salida. S-NSGA-II construye, además, una máscara de densidad organizada en franjas contiguas del vector de decisión, capaz de apagar de una sola vez filas completas de $W_1$, es decir, neuronas ocultas enteras; su operador de mutación de dispersión muta el nivel de dispersión objetivo de todo el individuo y luego apaga o reactiva pesos al azar hasta alcanzar ese nuevo nivel, otra vez sin mirar a qué neurona pertenece cada uno; y su operador de cruce intercambia al azar, entre los dos hijos, el estado activo o inactivo de los genes en los que los padres difieren. MOEA-CKF, por su parte, construye la máscara binaria de su primera generación activando un número aleatorio $k \in [1, D]$ de variables sobre el total de la red —pudiendo dejar tan solo una sinapsis activa de las $D$ totales—, sin ninguna garantía de que ese reducido conjunto toque siquiera una neurona de salida.

La huella que deja cada algoritmo en los resultados agregados es consistente con su mecanismo particular. La máscara binaria de MOEA-CKF, gobernada por un parámetro $k$ global e indiferente al segmento del vector que activa, deja proporcionalmente más soluciones con alguna clase de salida completa sin ninguna conexión (56.5\% frente al 37.5\% de S-NSGA-II); la máscara en franjas contiguas de S-NSGA-II, en cambio, apaga con mayor frecuencia neuronas ocultas completas (87.4\% muertas en promedio, frente al 74.4\% de MOEA-CKF). Que dos mecanismos de dispersión distintos produzcan dos huellas de falla igualmente distintas es un indicio adicional de que el patrón observado proviene del diseño particular de cada operador, y no de una coincidencia compartida entre ambos algoritmos.

Una parte del fenómeno es, además, mecánica y esperable: dado que la complejidad estructural ($f_1$) es un objetivo minimizado explícitamente, el punto de menor complejidad de cada frente tiende a ubicarse en el extremo de máxima dispersión, y de hecho el 100\% de los frentes conserva su solución más dispersa con al menos una salida desconectada. Lo que sí resulta relevante para la tesis es que, en al menos dos de las ocho combinaciones de dataset y algoritmo (\texttt{ds2}$\times$SNSGAII y \texttt{ds3}$\times$MOEACKF), el fenómeno no está confinado al extremo del frente: la solución de mejor accuracy de cada corrida —la que en la práctica se desplegaría— presenta de forma sistemática una clase de salida sin ninguna conexión de entrada, es decir, una red que jamás puede predecir una de las dos clases sin importar la entrada que reciba.

La asimetría por clase observada en \texttt{ds2} y \texttt{ds3} —donde es sistemáticamente la misma clase la que queda desconectada en la mayoría de las soluciones, aunque en sentido opuesto entre ambos datasets— es consistente con que ambos presentan clases desbalanceadas: dado que ninguno de los operadores de dispersión descritos distingue entre clases al apagar pesos, la clase minoritaria de cada dataset queda estructuralmente más expuesta a perder toda su conexión de entrada.

\subsection{Trabajo futuro recomendado}

Ninguna de estas causas se corrige evaluando de mejor manera la red resultante, sino interviniendo directamente en los operadores que deciden qué sinapsis permanecen activas. Queda como trabajo futuro —no implementado ni evaluado en el marco de esta tesis— explorar al menos cuatro direcciones de modificación sobre la codificación y los operadores actuales de ambos algoritmos. La primera consiste en incorporar un operador de reparación que se aplique después de cada mutación o cruce y antes de evaluar al individuo, capaz de detectar neuronas muertas o salidas con grado de entrada nulo y reactivar determinísticamente al menos una conexión hacia cada una de ellas. La segunda apunta a una inicialización consciente de la conectividad, que garantice, al construir cada individuo inicial, al menos un camino activo completo entre la entrada y cada clase de salida, en lugar de apagar genes de forma uniforme o por franjas sin considerar la estructura resultante. La tercera dirección consiste en operadores de mutación y cruce que protejan explícitamente la última conexión activa hacia una neurona de salida, o hacia una neurona oculta que constituya el único puente hacia esa salida, impidiendo que se apague por azar. La cuarta, finalmente, propone incorporar la conectividad de entrada a salida directamente en la función de aptitud —ya sea como objetivo adicional o como penalización— o como restricción de factibilidad, en lugar de dejar que dicha conectividad emerja únicamente como efecto colateral de minimizar $f_1$.
